% ==========================================================================
% GENERATED by scripts/split_paper.py. Do not edit this file.
%
% Source: imported_papers/Neural_Computation__CNA/NECO-D-24-00304-Source.tex
% Part 16: Quantifying Noise Filtration
%
% Edits belong upstream, in the published paper. To change what the
% chapter includes, change imports/cna.toml and re-run
% the splitter.
% ==========================================================================
\section{Quantifying Noise Filtration}\label{sec:quant_noise_filtering}

Our evaluations focus on the efficacy of net fragments in suppressing added Gaussian noise and complementing missing or occluded patterns. For both evaluation experiments, we report {\it recall} and {\it precision}: Recall is defined as the proportion of neurons that maintain activity (on) when noise is added, and precision quantifies the extent to which neurons activated in the presence of noise were also active without noise.
To assess the system's noise reduction capability in the case of Gaussian noise, we additionally measure the percentage of flipped neurons that revert to their original activation state, referring to this measurement as the {\it noise reduction rate}.
To assess subtractive noise, we measure the similarity between the original (created without noise) and reconstructed (created with noise) feature activations at the spatial locations where pixels have been removed. We call this metric the {\it feature reconstruction rate}.


\paragraph{Filtering Gaussian Noise.}
\begin{figure}[ht!]
  \begin{center}
    \includegraphics[width=\textwidth]{NECO-D-24-00304-Figure.10}
  \end{center}
  \caption{\textbf{System performance as function of the noise level}. First row: noise reduction rate; second row: recall; third row: precision. Columns differ in the attenuation coefficient $\gamma$. The colors within each plot represent different activation biases $b^{(S2)}$. The dotted red line is the autoencoder, serving as baseline.}\label{fig:noise_results}
\end{figure}
Figure \ref{fig:noise_results} displays noise reduction rates, recall, and precision across different parameter configurations for different levels of added noise.
The metrics demonstrate that increasing the attenuation coefficient or the activation bias results in enhanced noise filtration.

The overall efficacy of noise reduction remains consistently high across most parameter configurations, exhibiting a constant noise reduction rate of $\geq \approx 95\%$. This rate persists regardless of the magnitude of the introduced noise. As the noise suppression rate is approximately constant, this means that if there is more noise in the input, there will also be more noise in the output, albeit limited to around 5\% of the input noise.

Furthermore, all CNA models (except the one with the lowest power factor and bias) outperform the autoencoder, demonstrating that they achieve performance that is competitive with state-of-the-art deep learning models.

\paragraph{Reconstructing Subtractive Noise.}
\begin{figure}[ht!]
  \begin{center}
    \includegraphics[width=\textwidth]{NECO-D-24-00304-Figure.11}
  \end{center}
  \caption{\textbf{Reconstruction of partial occlusion}. {\it First row:} feature reconstruction rate; {\it second row:} recall; {\it third row:} precision. Columns are distinguished by the attenuation coefficient $\gamma$, and line colors within plots distinguish activation bias $b^{(S2)}$. The dotted red line is the autoencoder model.}\label{fig:li_results}
\end{figure}

Figure \ref{fig:li_results} displays the feature reconstruction rate, precision, and recall for removed line fragments of different lengths. Having a lower activation bias or attenuation coefficient $\gamma$ leads to a better reconstruction of removed patterns, as these settings allow neurons with less support to activate.

Compared to the autoencoder baseline, the CNA model achieves higher recall and precision but a lower feature reconstruction rate.
This indicates that the autoencoder is superior at reconstructing the pixels that have been removed but that, in general, much more noise (other neurons are activated/deactivated) is introduced.

\paragraph{Interpreting Metrics.}
While the reported values for precision and recall may appear low, we argue that these fragments are of good quality and point out that precision and recall metrics have to be interpreted differently than in a classification context: 
A net fragment consisting of many neurons characterizes (a part of) an object and maintains this characterization even when many of its neurons are inversed (and precision and recall are low). This aligns with the findings of \citet{ahmad_properties_2015}, highlighting that binary distributed activations demonstrate high robustness, retaining accurate interpretation even in the presence of numerous flipped neurons.

In the case of the depicted straight lines, the feature neuron activations of the two most similar lines overlap with $21.2\%$. Therefore, achieving precision and recall values surpassing this threshold permits reliable discrimination even among similar lines, which is the case for the CNA but not for the autoencoder (e.g., see Figure \ref{fig:noise_results}).

\clearpage
\newpage
